% Notes:
% \begin{itemize}
%     \item How to design dynamic tariffs?
%     \item Include a part for the COP approximation, how it varies over the year?
%     \item Is the district heating system run on a constant supply temperature or is it run depending on the outside temperature?
%     \item Are there other components included in the system?
%     \item It should depend on the inertia of the buildings, the supply temperature of the DH system
%     \item Extend a for a higher demand, as outlook
%     \item Mention in the overleaf document what could be done and will be done as extension?
% \end{itemize}

% Using 2024 data, the incumbent monthly tariff is compared with a mean-price-adjusted dynamic tariff and an unconstrained PTES benchmark. Dynamic pricing [increases PTES utilisation from X to Y cycles], [changes operating cost by X%], and alters the interaction between direct heat imports, storage, and heat pumps. Technical constraints [quantified effect]. The findings show whether the marginal-cost signal creates sufficient price spreads to incentivise economically valuable and technically feasible PTES operation.

% Another version

% \begin{abstract}
% Pit thermal energy storage (PTES) can provide district-heating flexibility by shifting heat purchases across time, but its technical potential is not necessarily activated under existing tariffs. In Høje Taastrup, observed PTES utilisation is approximately 10 equivalent full cycles per year, compared with an operational ambition of 25--30 cycles. The incumbent monthly tariff provides no within-month incentive to shift imports between low- and high-cost hours, while storage losses further reduce the attractiveness of cycling. A dynamic tariff based on marginal heat-generation costs may provide the missing incentive by rewarding charging when external heat is inexpensive and discharging when it is expensive.

% This paper evaluates how dynamic marginal heat pricing affects the incentive, utilisation, and operational value of PTES. An hourly mixed-integer model of the Høje Taastrup district-heating system jointly optimises heat imports, a 3.3~GWh PTES, and two heat pumps. Heat demand, eastern Danish electricity prices, and upstream marginal heat-generation costs are soft-linked as external time series. Storage losses and realistic capacity, ramping, operating-mode, and switching constraints are represented explicitly.

% \textbf{Include another paragraph about the results}
% \end{abstract}



%%%%%%%%%%%%%%%%%%%%%%%%%%%%%%%%%%%%%%%%%%%%%%%%%%%%%%%%%
% Notes regarding the scientific gap and how to point to it
%%%%%%%%%%%%%%%%%%%%%%%%%%%%%%%%%%%%%%%%%%%%%%%%%%%%%%%%%

% a different possible version for the abstract

\begin{comment}
Pit thermal energy storage (PTES) can provide district-heating flexibility by shifting heat production relative to heat demand. In systems with combined heat and power generation, heat can be produced and stored when electricity prices are high and co-produced electricity is most valuable. In multi-operator district heating networks, however, system dispatch and commercial settlement may operate at different temporal resolutions. Periodic cost-recovery tariffs may therefore not transmit the time-varying value of storage to the operator controlling its use.

This coordination problem is examined in Greater Copenhagen, where Varmelast coordinates system-wide heat production, VEKS charges the Høje Taastrup PTES based on this dispatch, and Høje Taastrup District Heating schedules its discharge under a monthly heat tariff. This paper develops a dynamic heat tariff based on a marginal-cost proxy for upstream heat production, accounting for fuel and operating costs, carbon costs, conversion efficiencies and the market value of co-produced electricity.

The tariff is soft-linked to an hourly mixed-integer optimisation model of the Høje Taastrup district heating system. The model jointly optimises direct heat imports, a 3.3 GWh PTES and two heat pumps under realistic operating constraints using heat-demand and electricity-price data for 2024. Monthly and dynamic tariff structures are compared using level-matched scenarios, while their effects on local expenditure, upstream revenue and a common operational-cost metric are evaluated separately.
\end{comment}


% Notes moved out of the compiled discussion during article formatting.

what is VEKS paying for the heat

looks suboptimal to just operate the heat pumps by default

heres space for further work now



% We need information on the internal constraints of the HTF network,
% particularly the achievable PTES discharge rate and the factors that
% influence it. What would be required to reach a discharge rate of
% 30~MW, and why has the maximum observed discharge rate been limited
% to approximately 23~MW?

% We also need data on:
% \begin{itemize}
%     \item Forecast heat demand in the HTF network.
%     \item The resulting PTES charging and discharging schedules.
%     \item Actual PTES charging and discharging.
% \end{itemize}

% A consistent dataset covering one month to one year would be highly
% valuable for developing and evaluating the model.

% \subsubsection{Distinguishing costs, storage value and tariffs}

% Three different quantities should be distinguished in the system model:

% \begin{center}
% \begin{tabularx}{\textwidth}{
%     @{}
%     >{\raggedright\arraybackslash}p{0.35\textwidth}
%     >{\raggedright\arraybackslash}X
%     @{}
% }
% \toprule
% \textbf{Quantity} & \textbf{Meaning} \\
% \midrule
% Current marginal heat-production cost
% & Cost of supplying an additional unit of heat now. \\[0.5em]

% Marginal value of stored heat
% & Benefit of retaining an additional unit of heat for future use. \\[0.5em]

% HTF discharge tariff
% & What HTF pays when receiving heat from the PTES. \\
% \bottomrule
% \end{tabularx}
% \end{center}

% The value of stored heat depends on expected future costs, demand,
% storage availability and operating restrictions. Varmelast's
% description explicitly links the PTES weekly outlook to this
% forward-looking valuation.

% \subsubsection{Understanding conservative discharge operation}

% It is essential to establish with HTF why discharge is operated
% conservatively. This is a prerequisite for further model development.
% If technical constraints prevent the higher discharge rates required
% by the proposed approach, the project should not proceed in its
% current form.

% Different reasons for conservative operation require different
% representations in the model:

% \begin{center}
% \begin{tabularx}{\textwidth}{
%     @{}
%     >{\raggedright\arraybackslash}X
%     >{\raggedright\arraybackslash}X
%     @{}
% }
% \toprule
% \textbf{Reason for conservative operation}
% & \textbf{Appropriate representation} \\
% \midrule
% A discharge rate would violate a network or temperature limit.
% & Hard physical constraint. \\[0.5em]

% Operators prefer a lower rate because they are uncertain about
% operating conditions.
% & Perceived risk or a preferred operating range. \\[0.5em]

% Operators avoid frequent or rapid adjustments.
% & Adjustment cost, deadband or limited decision frequency. \\[0.5em]

% Operators prefer an established discharge procedure.
% & A reference operating policy with a cost of departing from it. \\
% \bottomrule
% \end{tabularx}
% \end{center}

% \noindent
% \textbf{A financial penalty can change a preference. It cannot remove
% a physical constraint.}

% We should therefore ask HTF what ``conservative'' means operationally:
% lower discharge rates, slower ramps, earlier shutdown, larger safety
% margins or fewer interventions. These behaviours imply different
% model structures.

% Three distinct concepts should be described:

% \begin{center}
% \begin{tabularx}{\textwidth}{
%     @{}
%     >{\raggedright\arraybackslash}p{0.20\textwidth}
%     >{\raggedright\arraybackslash}X
%     >{\raggedright\arraybackslash}X
%     @{}
% }
% \toprule
% \textbf{Concept}
% & \textbf{Meaning in this case}
% & \textbf{Mathematical representation} \\
% \midrule
% Physical flexibility
% & Which discharge trajectories HTF can technically implement.
% & Feasible set defined by demand, temperature, storage and ramp
% constraints. \\[0.5em]

% Operating preference
% & Which feasible trajectories HTF prefers.
% & Operating rule or additional objective terms. \\[0.5em]

% Responsiveness
% & How much HTF changes its preferred operation when incentives change.
% & Parameters governing the trade-off between preferences and payments. \\
% \bottomrule
% \end{tabularx}
% \end{center}

% \subsubsection{Initial operational hypothesis}

% An initial operational hypothesis should be developed with HTF.
% For example, the relevant behaviour might be:

% \begin{quote}
% We can discharge faster, but normally maintain a lower discharge rate
% because it gives us a larger operating margin.
% \end{quote}

% Alternatively, it might be:

% \begin{quote}
% The discharge level itself is not the problem. We avoid changing it
% frequently because every adjustment requires intervention.
% \end{quote}

% The following questions should therefore be clarified:
% \begin{itemize}
%     \item What exactly do operators avoid: high discharge rates,
%     rapid changes, frequent changes or uncertainty?
%     \item Does this preference depend on demand, temperatures or
%     network configuration?
% \end{itemize}

% \subsubsection{Tariff cases and comparison benchmarks}

% A simple compensation tariff should be included as a serious
% comparator. Moving directly from the monthly tariff to a sophisticated
% dynamic tariff raises the question of whether the additional complexity
% is necessary. A simple discharge discount should therefore also be
% investigated.

% \begin{center}
% \begin{tabularx}{\textwidth}{
%     @{}
%     >{\raggedright\arraybackslash}X
%     >{\raggedright\arraybackslash}X
%     @{}
% }
% \toprule
% \textbf{Case} & \textbf{Purpose} \\
% \midrule
% Existing arrangements
% & Establish the baseline. \\[0.5em]

% Pumping-cost compensation or a constant discharge discount
% & Test whether a simple payment resolves the incentive gap. \\[0.5em]

% Dynamic discharge discount
% & Identify the additional value of varying the incentive over time. \\[0.5em]

% Dynamic discount plus shortfall charge
% & Identify the additional effect of liability. \\[0.5em]

% Feasible schedule-following benchmark
% & Measure the remaining scope for improvement. \\
% \bottomrule
% \end{tabularx}
% \end{center}

% \subsubsection{Proposed soft-coupled modelling process}

% The full process of the proposed soft-coupled model would look
% as follows:

% \begin{center}
% \small
% \begin{tabularx}{\textwidth}{
%     @{}
%     >{\raggedright\arraybackslash}p{0.20\textwidth}
%     >{\raggedright\arraybackslash}X
%     >{\raggedright\arraybackslash}X
%     @{}
% }
% \toprule
% \textbf{Step}
% & \textbf{Model or calculation}
% & \textbf{Main output} \\
% \midrule
% \textbf{1. Central scheduling}
% & PyPSA optimises represented heat supply and PTES operation using
% forecasts and HTF's declared capability
% & Charging/discharging schedule and estimated storage value \\[0.7em]

% \textbf{2. Tariff calculation}
% & Translate the economic signals into a candidate discharge tariff
% and deviation charge
% & Prices and settlement rules announced to HTF \\[0.7em]

% \textbf{3. HTF response}
% & Behavioural model balances payments, operating preferences and
% local feasibility
% & Actual discharge, heat-exchanger supply and any adjustable local
% production \\[0.7em]

% \textbf{4. System adjustment}
% & Calculate how actual operation changes upstream supply and
% storage availability
% & Corrective operating costs, updated storage state and stakeholder
% outcomes \\[0.7em]

% \textbf{5. Tariff evaluation}
% & Repeat for alternative tariff parameters and operating conditions
% & Best-performing structures under the selected criteria \\
% \bottomrule
% \end{tabularx}
% \end{center}


decouple it from the big system and show that is has to be on a smaller side, with \\
mention constrained and unonstrianed model but it shouldnt be focus of the whole paper \\
lower heat pump cost in the distribution side, use less chp units on the transmission system \\ show the bigger argument
show on the economic impact of the system then

Energy -

seperating two problems one for the local systme and the global system
marginal heat price send from the bigger system to the small

incldue the cost of the pump for the PTES.


discount on the extra heat - inlcude pump cost into the system - or ramp cost for the reward - if the dual variables are useful for the constraints or heat exchanger, how to explore this then. price or reward for the pumps to operate

put discussions in the model and see how it changes.
